% from ega4-1-fr.tex
\oldpage[IV-1]{243}

où $X=\operatorname{Spec}(A)$ et où $E=f(X')$, où
$X'=\operatorname{Spec}(A')$ est affine, $f:X'\to X$ un morphisme
correspondant à un homomorphisme d'anneaux $A\to A'$. Notons maintenant que
$A'$ est limite inductive de la famille $(A'_\lambda)$ de ses
sous-$A$-algèbres \emph{de type fini}, ordonnée par inclusion~; si
$X'_\lambda=\operatorname{Spec}(A'_\lambda)$, $f$ se factorise en
$X'\to X'_\lambda\xrightarrow{f_\lambda}X$ et il résulte de (1.9.2, (ii))
que l'on a $f(X')=\bigcap_\lambda f_\lambda(X'_\lambda)$~; si l'on démontre
que $f_\lambda(X'_\lambda)$ est intersection de parties constructibles de
$X$, il en sera de même de $f(X')$. On peut donc se borner au cas où $A'$ est
une $A$-algèbre \emph{de type fini}. Or, on a le lemme suivant~:
\end{proof}

\begin{lemma}[1.9.3.1]
\label{IV.1.9.3.1-fr}
Soient $A$ un anneau, $A'$ une $A$-algèbre de type fini. Alors $A'$ est
$A$-isomorphe à une limite inductive filtrante $\varinjlim A''_\mu$ où les
$A''_\mu$ sont des $A$-algèbres de présentation finie (1.4.1).
\end{lemma}

\begin{proof}
En effet, on peut écrire $A'=B/\mathfrak J$, avec
$B=A[T_1,\ldots,T_n]$ et $\mathfrak J$ un idéal de $B$. Mais $\mathfrak J$
est limite inductive des idéaux \emph{de type fini} $\mathfrak J'_\mu$ de
$B$, contenus dans $\mathfrak J$, ordonnés par inclusion~; comme dans la
catégorie des $B$-modules le foncteur $\varinjlim$ est exact, on en conclut
que $A'=\varinjlim B/\mathfrak J'_\mu$ à un $A$-isomorphisme près.
\end{proof}

Le même raisonnement que ci-dessus permet alors de se ramener au cas où $A'$
est une $A$-algèbre de présentation finie~; mais alors $f(X')$ est
constructible dans $X$ en vertu du théorème de Chevalley (1.8.5), ce qui
prouve $b)$.

Démontrons enfin que $b)$ entraîne $a'')$. Si $E$ est intersection d'une
famille $(G_\lambda)_{\lambda\in L}$ de parties constructibles de $X$,
chaque $E_\alpha$ est intersection des $G_\lambda\cap X_\alpha$, qui sont
constructibles dans $X_\alpha$ (0\textsubscript{III}, 9.1.8, (i)), donc on
est ramené au cas où $X=\operatorname{Spec}(A)$ est affine.

On sait alors (1.8.3) que pour tout $\lambda\in L$ il existe un morphisme
$f_\lambda:X'_\lambda\to X$, où
$X'_\lambda=\operatorname{Spec}(A'_\lambda)$ est affine, tel que
$f_\lambda(X'_\lambda)=G_\lambda$. Il suffit donc de prouver le lemme
suivant~:

\begin{lemma}[1.9.3.2]
\label{IV.1.9.3.2-fr}
Soient $(A'_\lambda)_{\lambda\in L}$ une famille de $A$-algèbres,
$X=\operatorname{Spec}(A)$, $X'_\lambda=\operatorname{Spec}(A'_\lambda)$,
et soit $f_\lambda:X'_\lambda\to X$ le morphisme structural. Pour toute
partie finie $J$ de $L$, posons
$A'_J=\bigotimes_{\lambda\in J}A'_\lambda$ (produit tensoriel de
$A$-algèbres), $X'_J=\operatorname{Spec}(A'_J)$, et soit $f_J:X'_J\to X$ le
morphisme structural. On a alors
$f_J(X'_J)=\bigcap_{\lambda\in J}f_\lambda(X'_\lambda)$. Si
$A'=\varinjlim A'_J$, $J$ parcourant l'ensemble filtrant des parties finies
de $L$, $X'=\operatorname{Spec}(A')$, et si $f:X'\to X$ est le morphisme
structural, on a $f(X')=\bigcap_{\lambda\in L}f_\lambda(X'_\lambda)$.
\end{lemma}

\begin{proof}
La première assertion résulte de (I, 3.4.7)~; la deuxième résulte de (1.9.2,
(ii)), qui donne la relation $f(X')=\bigcap_Jf_J(X'_J)$.
\end{proof}

\begin{definition}[1.9.4]
\label{IV.1.9.4-fr}
Soit $X$ un espace topologique. On dit qu'une partie $E$ de $X$ est
\emph{pro-constructible} (resp. \emph{ind-constructible}) dans $X$ si, pour
tout $x\in X$, il existe un voisinage ouvert $U$ de $x$ dans $X$ tel que
$E\cap U$ soit intersection (resp. réunion) de parties localement
constructibles de $U$.
\end{definition}

Tenant compte de (1.2.7) et de (0\textsubscript{III}, 9.1.8 et 9.1.12), les
conditions équivalentes de la proposition (1.9.3) s'expriment donc en disant
que $E$ est \emph{pro-constructible} dans $X$, les ensembles localement
constructibles dans $X$ étant identiques aux ensembles constructibles dans
$X$ sous les hypothèses de (1.9.3).

